\documentclass[10pt,twocolumn]{article}
\usepackage[utf8]{inputenc}
\usepackage[T1]{fontenc}
\usepackage{lmodern}
\usepackage[a4paper,margin=0.68in]{geometry}
\usepackage{amsmath,amssymb,amsthm,mathtools}
\usepackage{graphicx}
\usepackage{authblk}
\usepackage{booktabs}
\usepackage[font=small,labelfont=bf]{caption}
\usepackage{tabularx}
\usepackage{enumitem}
\usepackage{stfloats}
\usepackage{microtype}
\usepackage{xcolor}
\usepackage[round,authoryear]{natbib}
\bibpunct{(}{)}{;}{a}{,}{,}
\usepackage[colorlinks=true,linkcolor={blue!55!black},citecolor={green!45!black},urlcolor={blue!60!black}]{hyperref}

\renewcommand{\arraystretch}{1.15}
\renewcommand{\topfraction}{0.92}
\renewcommand{\bottomfraction}{0.75}
\renewcommand{\textfraction}{0.06}
\renewcommand{\floatpagefraction}{0.8}
\renewcommand{\dbltopfraction}{0.92}
\renewcommand{\dblfloatpagefraction}{0.8}
\setcounter{topnumber}{3}
\setcounter{totalnumber}{5}
\setcounter{dbltopnumber}{3}

\DeclareMathOperator*{\argmax}{arg\,max}
\providecommand{\doi}[1]{\href{https://doi.org/#1}{\nolinkurl{doi:#1}}}
\newcommand{\E}{\mathbb{E}}
\newcommand{\filt}{\mathcal{F}}
\newcommand{\Nb}{\mathcal{N}}
\newcommand{\Pk}{P_k}

\newtheorem{theorem}{Theorem}
\newtheorem{proposition}{Proposition}
\newtheorem{corollary}{Corollary}
\theoremstyle{definition}
\newtheorem{assumption}{Assumption}
\newtheorem{definition}{Definition}

\title{\textbf{When Does Design Thinking Work?\\A Convergence Theorem for the Iterate--Test Core}}
\author[1]{Paweł Kuraś}
\author[2]{Adrian Michalski}
\author[3]{Alicja Gerka}
\author[3]{Patryk Organiściak\thanks{\texttt{org@prz.edu.pl}}}
\affil[1]{Institute of Telecommunications, AGH University of Science and Technology, Kraków, Poland}
\affil[2]{Faculty of Mathematics and Applied Sciences, Rzeszów University of Technology, Rzeszów, Poland}
\affil[3]{Department of Complex Systems, Rzeszów University of Technology, Rzeszów, Poland}
\date{}

\begin{document}
\maketitle


\begin{abstract}
\noindent Design thinking is taught as a loop of ideation, prototyping and testing and is credited with reliably improving what teams build, yet there is no account of the conditions under which such a loop must improve an outcome. We give such an account. Within a single framing of the problem we model the loop as a search over a finite space of prototypes whose utility is learned from noisy tests. The prototype a team believes to be best is eventually a fixed one, and a local optimum of utility, when five conditions hold: the frame is finite, ideation stays close to what has been built, the supply of improving ideas does not dry up too quickly, tests are unbiased and repeated, and nothing that has been tested is thrown away. Biased tests, and criteria that drift within bounds, weaken the result by a bounded amount, and with several stakeholders the limit is a locally Pareto-optimal set. Simulations on three classes of landscape, whose local optima are known exactly, show what the theorem leaves out. Except on very rugged landscapes the limit is rarely reached, even after thousands of rounds. At every horizon a team has already tested a better prototype than the one it would choose, a loss comparable to its failure to find better ones, and it overestimates the prototype it has chosen by 4 to 8\% of the optimum. Under a fixed budget of tests the best policy builds one to three new prototypes for every five tests and spends the rest re-testing the leading ones, and testing every prototype once is the worst policy. A bias towards the current favourite is harmless while it is smaller than the noise of a test and costly beyond that. Given the same tests, the loop far outperforms random search and hill-climbing and, when both screen their ideas, comes within one and a half percent of the optimum of a tuned bandit policy; a rough screening of ideas before building is worth more than the choice between the two. Reframing once improves the result by about one percent, and an expected-improvement rule for when to reframe does no better than a fixed plan. Whether design teams meet the five conditions is an empirical question, and we describe the measurements that would answer it.
\end{abstract}

\vspace{0.4em}
\noindent\textbf{Keywords:} design thinking; iteration; prototyping; Bayesian decision theory; convergence; rugged landscapes; exploration and exploitation

\section{Introduction}

Design thinking asks teams to understand the people they are designing for, frame a problem, generate many ideas, build cheap prototypes and test them with users, and to repeat the last three steps until something works \citep{dschool2010}. The method has spread from product design into management, public policy and health care \citep{liedtka2015, mintrom2016, altman2018}, and it has done so without a theory of why it should work. Reviews describe a construct that is fragmented and thinly grounded \citep{johansson2013, micheli2019}. In health care, where the evidence has been reviewed systematically, it consists mostly of small studies of fair or poor quality \citep{altman2018}. The process is presented as a sequence of steps and also as a set of overlapping spaces to be revisited as needed \citep{brown2008}, which makes it hard to say what would count as doing it wrong.

This paper asks a narrow question that can be answered exactly. \emph{Under what conditions does the iterate--test core of design thinking provably improve what a team builds, and what does the guarantee leave out?} We take the framing of the problem as given and study the loop that follows it: propose variants, build one, test it, revise what is believed about every prototype built so far, and repeat. A prototype is a set of features. Its value is a utility that the team does not know and can observe only through noisy tests, and the team keeps a belief about each prototype it has tested. In this setting improvement has a definite meaning, and so does a local optimum: a prototype that no nearby variant beats.

The first part of the paper is a convergence result. Suppose the team entertains finitely many features, generates ideas near what it has already built, keeps producing improving ideas at a rate that is not summable, tests without bias and re-tests what it has kept, and never discards a tested prototype. Then, with probability one, the prototype the team believes to be best is eventually a fixed one, and it is a local optimum (Theorem~\ref{thm:converge}). The argument uses three standard facts: in a finite space the set of prototypes stops growing, repeated tests make beliefs correct, and a conditional Borel--Cantelli argument excludes a limit that a better variant could still replace. Results of this kind are known for search that retains its best solution and for stochastic search under noise. We state the result for the design loop, where each condition names something a team does and each can be weakened to see what is lost: a bias of at most $B$ in the tests leaves the team at a prototype that no variant beats by more than $2B$; criteria that drift within a band of half-width $D$ behave like a bias of size $D$; and when stakeholders disagree, the set of non-dominated prototypes converges to a locally Pareto-optimal set.

The second part measures what the theorem leaves open, by simulation on three classes of landscape whose local optima are known exactly. The limit turns out to be remote. On landscapes built from desirability, feasibility and viability, only a small minority of runs reach a local optimum even after thousands of rounds, so what a team has when it stops depends on how fast the loop climbs. Three results follow. A team's choice among its own prototypes costs it about as much as its failure to find better ones, and on smooth landscapes more, while its estimate of the prototype it has chosen is systematically too high. Under a fixed budget of tests, never testing anything twice is the worst policy, and unless tests are very precise the best policies spend most of the budget re-testing the leading prototypes. A bias in favour of the current favourite does no measurable harm until it is about as large as the noise of a test. We also compare the loop with bandit and local-search policies given the same tests, and a team that never reframes with teams that do.

None of this shows that design teams satisfy the conditions. The model says what a team would have to do for iteration to be guaranteed to help, and Section~\ref{sec:empirical} sets out what would have to be measured in real projects to test it.

\section{Background and related work}
\label{sec:related}

Three bodies of work bear on the question: descriptions of design thinking as a process, formal theories of design, and the study of search and testing under uncertainty.

The five modes of the Stanford d.school (empathize, define, ideate, prototype, test) \citep{dschool2010} and the Design Council's Double Diamond, which alternates divergent and convergent phases across discovery, definition, development and delivery \citep{designcouncil2007}, are two widely taught descriptions. Brown adds the criteria by which an outcome is judged: what people want, what can be built, and what a business can sustain \citep{brown2008, brown2009}. An interview study of large firms \citep{carlgren2016} and a review of the literature \citep{micheli2019} both find experimentation and iteration among the elements that recur in accounts of the method. \citet{liedtka2015} proposes that it improves outcomes by mitigating known biases of judgement, among them the tendency to seek confirmation for a solution one already favours; designers shown an example solution are known to reproduce its features, flaws included \citep{jansson1991}. Experimental work bears on particular steps. Novice teams differ in how they combine need-finding, brainstorming and prototyping and in how they reflect on them \citep{seidel2013}; designers who built several prototypes before receiving feedback produced better final designs than those who refined one at a time \citep{dow2010}; and groups that generate more ideas do not thereby select better ones \citep{rietzschel2006}. Our model takes the last three modes as its object and treats the first two as producing the frame within which they operate.

Formal theories of design begin with Simon, whose science of the artificial \citep{simon1996} casts design as search under bounded rationality, ending in a satisfactory alternative \citep{simon1956}. Later accounts describe what design reasoning operates on. The function--behaviour--structure ontology separates what an artefact is for, how it behaves and what it is \citep{gero2004}; C--K theory describes design as the joint expansion of a space of concepts and a space of knowledge \citep{hatchuel2009}; axiomatic design gives criteria for a good mapping from requirements to parameters \citep{suh1990}; and models of design iteration identify which coupled tasks force rework \citep{smitheppinger1997}. Co-evolution models hold that problem and solution develop together \citep{maherpoon1996, dorstcross2001}. \citet{schon1983} places the setting of the problem at the centre of professional practice, \citet{dorst2011} describes framing as an abductive step, and experienced designers report reframing the briefs they are given as a matter of course \citep{paton2011}. Computational models of design teams exist as well; \citet{mccomb2015} simulate a team as a set of agents performing simulated annealing over a solution space and compare its behaviour with human data. These theories and models describe how designers reason and search. They do not state conditions under which an iterative process must improve its outcome. Our model is the simplest representation in which that question can be posed: design as exploration of a space of alternatives \citep{woodbury2006}, with the frame held fixed.

The convergence of a search process to a local optimum is well understood outside design. Adaptive walks on rugged landscapes end at local optima \citep{kauffmanlevin1987}, and in the NK model the optima multiply and the walks shorten as interactions among components increase \citep{kauffmanweinberger1989, kauffman1993}. Organisational search of this kind adapts well locally and misses distant peaks \citep{levinthal1997, levinthalmarch1993}, which is the exploitation side of March's trade-off \citep{march1991}. A canonical genetic algorithm does not converge to the global optimum, while variants that always retain the best solution found do \citep{rudolph1994}, and retention plays a comparable role in our theorem. In design a prototype's value is estimated, not observed. Search under noisy evaluation has its own literature: simulated annealing converges under noise if the noise is driven down fast enough, for example by repeated observation \citep{gutjahr1996}; random search over a discrete set converges when samples grow at a suitable rate \citep{homem2003}; evolutionary algorithms cope with noisy fitness by re-sampling candidates \citep{jin2005}; and in ranking and selection the efficient allocation of a simulation budget gives most replications to the leading alternative and its close competitors \citep{chen2000}. In organisational search, imperfect evaluation changes the character of the process and can even help it leave a local peak \citep{knudsen2007}. The economics of testing in product development have been studied in their own right: testing as a sequence of experiments whose mode can be chosen \citep{thomke1998}, testing alternatives in parallel or in sequence \citep{loch2001}, the number of concepts worth testing in parallel \citep{dahan2001}, and the way correlated designs shape a test sequence \citep{erat2008}. We prove the corresponding result for the design thinking loop and study the same questions by simulation.

A team that maintains beliefs about alternatives and chooses what to test next is also solving a sequential decision problem of the kind studied in Bayesian optimization \citep{jones1998, shahriari2016}, where expected improvement and the knowledge gradient \citep{frazier2008, frazier2009} are standard criteria and regret bounds are available for smooth objectives \citep{srinivas2012}. In a finite space with independent beliefs the closer analogy is best-arm identification, in which the number of tests needed to separate two alternatives grows, up to logarithmic factors, with the inverse square of the difference between them \citep{evendar2006, mannor2004, audibert2010}. We use these results to explain why choosing among prototypes is slow, and to supply the policies with which the loop is compared \citep{auer2002, russo2018}.

\section{The model}
\label{sec:model}

The model has two levels, the choice of a frame and the loop that runs inside it, and three ingredients: prototypes, a utility, and tests.

\subsection{Frames, prototypes and the loop}

A project moves through the five modes of design thinking, and it often moves backwards: a test that reveals a misread need sends the team back to its users, and one that reveals a weak concept sends it back to ideation (Figure~\ref{fig:process}). We split this movement into two levels. The outer level is the choice of a \emph{frame}. Empathizing gathers observations about users, and defining commits the team to a point of view, a statement of whose problem is being solved and what would count as solving it. In Dorst's account this step is abductive: starting from the value it wants to create, the team conjectures both the thing to be made and the working principle by which that thing would deliver the value \citep{dorst2011, kolko2010}. We do not model how the conjecture is produced. For our purposes its result is a pair $(F,U)$: a finite set $F$ of features that solutions may combine, and a utility $U$ that says how good each combination is. A \emph{reframe} replaces $(F,U)$ by another pair, and a \emph{framing epoch} is the stretch of the project between two reframes.

\begin{figure}[t]
\centering
\includegraphics[width=\columnwidth]{Fig1.png}
\caption{The design thinking process as modelled. Solid arrows are the forward sequence; dashed arrows are returns. The theorem concerns the loop among Ideate, Prototype and Test within one framing; a return to Empathize or Define starts a new framing epoch.}
\label{fig:process}
\end{figure}

The inner level is the loop inside an epoch. A \emph{prototype} is a subset $p\subseteq F$, so the space of prototypes is $2^F$. In round $k$ the team does three things. It ideates: it proposes \emph{ideas}, each a variant of a prototype it has already built, which is the divergent step. It prototypes: it chooses one of its ideas and builds it. It tests: it gathers evidence about the new prototype and about some of the older ones, revises its beliefs, and so decides which prototype is currently its \emph{favourite}, which is the convergent step.

Consider a team designing a travel mug that has settled on the commuter who drinks while walking, and whose frame admits four features: a folding handle, a one-touch lid, a spill seal and a temperature display. There are sixteen prototypes, from the bare mug to the mug with everything. If ideation changes one feature at a time, the variants of the mug with a seal are the bare mug and the three mugs that add one further feature to the seal. The mug with a lid and a seal is a local optimum if adding the handle or the display, or removing the lid or the seal, makes it worse on balance, even if the mug with a handle and a display and no lid, three changes away, would be better still. Deciding that the commuter is the wrong user is a reframe.

We write $\Pk\subseteq 2^F$ for the set of prototypes built and tested by the end of round $k$, and $\filt_k$ for the history up to that point: the ideas proposed, any rough estimates made of them, the prototypes built, and the tests and their results. Table~\ref{tab:notation} collects the notation.

\begin{table}[t]
\caption{Notation.}
\label{tab:notation}
\small
\begin{tabularx}{\columnwidth}{@{}lX@{}}
\toprule
$F$, $2^F$ & features admitted by the frame; prototype space\\
$U(p)$ & utility of prototype $p$, in $[0,1]$; $U^*=\max_p U(p)$\\
$\Nb(p)$ & variants of $p$ that ideation can reach in one step\\
$\Pk$, $P_\infty$ & prototypes built and tested by round $k$; their eventual set\\
$\filt_k$ & history up to the end of round $k$\\
$M^p_k$ & belief about $p$ after round $k$ ($\widehat M^p_k$ under biased tests)\\
$p^*_k$, $X_k$ & the favourite, $\argmax_{p\in \Pk}M^p_k$, and the belief about it\\
$\delta_k$ & lower bound on the chance that round $k$ builds a better variant of the favourite\\
$\sigma$ & standard deviation of the noise of a test\\
$b_k$, $B$ & bias of the test in round $k$ and its bound; $b$ is the constant bias used in simulations\\
\bottomrule
\end{tabularx}
\end{table}

\subsection{Utility}

Following \citet{brown2009}, a solution is judged by its desirability to users, its technical feasibility and its business viability. In our examples the utility is a weighted sum of the three,
\begin{equation}
\label{eq:utility}
U(p)=w_{\mathrm d}U_{\mathrm d}(p)+w_{\mathrm f}U_{\mathrm f}(p)+w_{\mathrm v}U_{\mathrm v}(p),
\end{equation}
with each component in $[0,1]$ and weights summing to one, as in multi-attribute utility theory \citep{keeney1976}. The theory does not depend on this form. It uses two properties of $U$ only: it is bounded, and it does not change during an epoch. A utility in which a failure on one criterion cannot be compensated by the others has both. Section~\ref{sec:pareto} drops the assumption that stakeholders agree on a single $U$.

\subsection{Tests and beliefs}

The team does not know $U$. A test of prototype $p$ in round $k$ returns
\begin{equation}
\label{eq:test}
y=U(p)+b_k+\varepsilon_k ,
\end{equation}
where $\varepsilon_k$ is noise that has mean zero and variance at most $\sigma^2$ given the history and the choice of what to test, and $b_k$ is a systematic error, possibly depending on the prototype and on the history, with $|b_k|\le B$. A user interview, an expert critique and a feasibility estimate are all tests in this sense. They differ in $\sigma$ and in how prone they are to bias.

From the results of its tests the team forms a belief $M^p_k$ about the utility of each prototype it has built: an estimate, such as the average of the results or the posterior mean of a Bayesian model \citep{degroot2004}. Its favourite $p^*_k$ is the prototype with the highest belief, ties being broken by any fixed rule, and
\begin{equation}
\label{eq:Xk}
X_k=\max_{p\in \Pk}M^p_k
\end{equation}
is how good the team thinks its favourite is. Two other quantities must be kept apart from $X_k$. The true utility $U(p^*_k)$ of the favourite is what the team would deliver if it stopped now, and we call it the \emph{shipped} utility. The best true utility among all tested prototypes, $\max_{p\in\Pk}U(p)$, is what the team has \emph{found}, whether or not it knows it.

An idea that has not been built has no test results. The choice of which idea to build next must therefore rest on something cheaper, such as a rough estimate, an analogy or a critique, and we call this choice \emph{pre-screening}. It affects the speed of the loop and not its limit.

\section{Convergence within a frame}
\label{sec:theory}

This section states the conditions as five assumptions, proves the result, and examines what remains of it when the assumptions are weakened.

\subsection{Five assumptions}

\begin{assumption}[Finite frame]
\label{as:finite}
Within an epoch the feature set $F$ is finite, the utility $U:2^F\to[0,1]$ does not change, and no two prototypes have exactly the same utility.
\end{assumption}

The exclusion of ties is a convenience; it holds for almost every utility when criteria are measured on a continuous scale. Finiteness is a real restriction, because it rules out the arrival of features nobody had thought of. If features are added during an epoch but the additions stop, the results apply from the last addition onward. If they never stop, there is no fixed space in which to speak of an optimum.

\begin{assumption}[Local ideation]
\label{as:local}
Each prototype $p$ has a set $\Nb(p)\subseteq 2^F$ of variants, and every idea, apart from the ideas a team starts with, is a variant of a prototype already in $\Pk$. A prototype is \emph{locally optimal} if none of its variants has higher utility.
\end{assumption}

In the simulations the variants of $p$ are the prototypes that differ from it in at most two features. Locality is what makes the limit a local optimum. If every prototype were a variant of every other, the same argument would give convergence to the global optimum, at a price in speed that Section~\ref{sec:ablation} measures.

\begin{assumption}[Persistent generativity]
\label{as:gen}
There are numbers $\delta_k\ge0$, determined by the history $\filt_{k-1}$, such that whenever the favourite $p^*_{k-1}$ has a better variant that has not yet been built, round $k$ builds and tests such a variant with conditional probability at least $\delta_k$ given $\filt_{k-1}$, and
\begin{equation}
\label{eq:nonsummable}
\sum_{k\ge1}\delta_k=\infty .
\end{equation}
\end{assumption}

The assumption allows ideas to become harder to find as the team closes in on a peak. A chance of improvement that falls like $1/k$ satisfies \eqref{eq:nonsummable}; one that falls like $1/k^2$ does not. A constant bound $\delta_k\ge\delta>0$ is the special case in which returns do not diminish at all. The assumption concerns building only. The team need not recognise the better variant when it builds it, and need not know that one exists; recognition is the business of the next assumption.

\begin{assumption}[Unbiased, repeated testing]
\label{as:test}
Tests are unbiased ($B=0$), every prototype in $\Pk$ continues to be tested, so that each is tested infinitely often, and the belief about a prototype differs from the average of its test results by an amount that vanishes as its tests accumulate.
\end{assumption}

The last clause holds for the average itself and for the posterior mean of the usual conjugate models under any prior, since the weight of the prior vanishes. The consequence used below is that beliefs become correct. The noise accumulated over the tests of a given prototype is a martingale with increments of bounded variance, so by the strong law for martingales \citep[Sec.~4.5]{durrett2019} the average of the results converges to $U(p)$, and therefore $M^p_k\to U(p)$ almost surely for every prototype that is retained. Tests may be chosen adaptively, in any order that the history suggests.

\begin{assumption}[Retention]
\label{as:retention}
No tested prototype is discarded: $\Pk\subseteq P_{k+1}$.
\end{assumption}

Retention keeps the best prototype found so far available for comparison, which is the role that elitism plays in evolutionary search \citep{rudolph1994}.

\subsection{The theorem}

\begin{theorem}
\label{thm:converge}
Let Assumptions~\ref{as:finite}--\ref{as:retention} hold. Then almost surely the favourite is eventually a single prototype $p_\infty$, this prototype is locally optimal, and the team's belief about it converges to its utility, $X_k\to U(p_\infty)$.

If, in addition, $U$ is drawn from a prior distribution on utilities with values in $[0,1]$ and the beliefs $M^p_k=\E[U(p)\mid\filt_k]$ are the posterior means under that prior, then $X_k$ is a bounded submartingale with respect to the prior: the team never expects its belief about its favourite to fall.
\end{theorem}

\begin{proof}
By retention, $\Pk$ is a non-decreasing sequence of subsets of the finite set $2^F$, so there is a random round $k_0$ after which it equals a fixed set $P_\infty$. Every member of $P_\infty$ is tested infinitely often, so $M^p_k\to U(p)$ for each. The utilities of the members of $P_\infty$ are distinct, so from some round on the favourite is the member $p_\infty$ with the highest utility, and $X_k\to U(p_\infty)$.

Suppose that, with positive probability, $p_\infty$ has a variant $q$ with $U(q)>U(p_\infty)$. Such a $q$ cannot belong to $P_\infty$, of which $p_\infty$ is the best member. On this event the favourite has a better unbuilt variant in every late round, so by Assumption~\ref{as:gen} the conditional probabilities of building one sum to infinity. By L\'evy's extension of the Borel--Cantelli lemma \citep[Sec.~4.3]{durrett2019} one is then built after round $k_0$, which contradicts $\Pk=P_\infty$ for $k\ge k_0$. Hence $p_\infty$ is locally optimal almost surely.

For the second statement, $M^p_k$ is for each $p$ a martingale with values in $[0,1]$, and the set over which the maximum in \eqref{eq:Xk} is taken does not shrink, so by the conditional Jensen inequality
\begin{align*}
\E[X_{k+1}\mid\filt_k]&\ge\E\Big[\max_{p\in \Pk}M^p_{k+1}\,\Big|\,\filt_k\Big]\\
&\ge\max_{p\in \Pk}\E[M^p_{k+1}\mid\filt_k]=X_k .\qedhere
\end{align*}
\end{proof}

Each assumption does one job in the proof. Retention and finiteness make the set of prototypes settle. Repeated unbiased testing makes beliefs correct; without repetition a lucky first test could keep a mediocre prototype on top for ever. Generativity prevents the process from stopping short of a local optimum, and locality determines which optimum is meant.

The two statements of the theorem differ in kind. The first holds for every utility without ties and whatever the team's prior, because it uses only the fact that beliefs become correct. The second is the familiar fact that information has non-negative expected value to a Bayesian \citep{degroot2004}, and it holds only from the standpoint of a team whose prior is right and whose beliefs use everything in the history. A team whose prior is much more diffuse than the real spread of utilities overestimates its favourite, because the favourite is chosen for having the highest estimate and a diffuse prior does nothing to discount the luck in that estimate. The belief about any given favourite then falls when it is tested again. This is the optimizer's curse \citep{smithwinkler2006}, and Section~\ref{sec:dynamics} shows that it is large. The simulated team is of this kind: its prior is diffuse, its beliefs are normal posterior means confined to the unit interval, and it ignores its rough estimates once a prototype is built.

\subsection{Biased tests and drifting criteria}
\label{sec:bias}

Real tests are not unbiased. People seek and read evidence in ways that favour what they already believe \citep{nickerson1998}, and a team with a favourite is no exception. Suppose the bias in \eqref{eq:test} is bounded by $B$ and the team is unaware of it, so that it forms its beliefs as if tests were unbiased. Write $\widehat M^p_k$ for the resulting beliefs.

\begin{proposition}
\label{prop:bias}
Let Assumptions~\ref{as:finite}, \ref{as:local} and~\ref{as:retention} hold, let every retained prototype be tested infinitely often with $|b_k|\le B$, let beliefs be averages of test results or conjugate normal posterior means, and let Assumption~\ref{as:gen} hold with $\delta_k\ge\delta$ for a constant $\delta>0$. Then almost surely, from some round on, every favourite is $2B$-locally optimal: none of its variants has a utility exceeding its own by more than $2B$.
\end{proposition}

\begin{proof}
A belief is a weighted average of the prior mean and the test results, with the weight of the prior tending to zero. Shifting each result by at most $B$ shifts the average by at most $B$, and the unbiased part converges to $U(p)$ as before, so $\limsup_k|\widehat M^p_k-U(p)|\le B$ for every $p\in P_\infty$. Let $u$ be the highest utility in $P_\infty$. A member with utility below $u-2B$ has beliefs eventually below those of the best member, so every late favourite has utility at least $u-2B$. The favourite may keep changing among such members, which is why a constant $\delta$ is assumed: let $H_k$ be the event that the favourite at round $k-1$ has a better variant not yet built. Each time $H_k$ occurs a new prototype is built with probability at least $\delta$, so if $H_k$ occurred infinitely often, infinitely many prototypes would be built, by the Borel--Cantelli argument above, which is impossible in a finite space. Eventually, then, no favourite has a better unbuilt variant, and every variant that has been built has utility at most $u$, which exceeds the favourite's by at most $2B$.
\end{proof}

For $B=0$ the proposition gives the local optimality asserted in Theorem~\ref{thm:converge}. The bound is a worst case that lets the bias be arranged as unfavourably as possible, always against the better of two close prototypes. A bias that merely flatters the favourite is harmful only when it is large enough to overturn a comparison, and the simulated loss in Section~\ref{sec:ablation} is accordingly flat for small biases and steep for large ones.

The same argument covers criteria that move. Suppose that during an epoch the utility is $U_k$ at round $k$, for instance because the weights in \eqref{eq:utility} are renegotiated, and that it stays within a band around some reference utility $\bar U$ without ties.

\begin{corollary}
\label{cor:drift}
If $|U_k(p)-\bar U(p)|\le D$ for all $k$ and $p$, then under the conditions of Proposition~\ref{prop:bias}, with unbiased tests of the current utility and with ``better variant'' in Assumption~\ref{as:gen} understood with respect to $\bar U$, every late favourite is $2D$-locally optimal for $\bar U$ and $4D$-locally optimal for every $U_k$.
\end{corollary}

\begin{proof}
A test of $U_k(p)$ is a test of $\bar U(p)$ with a bias $U_k(p)-\bar U(p)$ of at most $D$, so Proposition~\ref{prop:bias} applies with $B=D$. For a late favourite $p$ and a variant $q$, $U_k(q)-U_k(p)\le\bar U(q)-\bar U(p)+2D\le4D$.
\end{proof}

This is as far as the model goes towards the co-evolution of problem and solution \citep{dorstcross2001}: it tolerates a problem that shifts within bounds while the solution is developed. Criteria that drift without bound are outside it, because a team whose notion of a good outcome keeps moving must discount old evidence, and a loop that discounts evidence is a different process. A discontinuous change of criteria is a reframe.

\subsection{Several stakeholders}
\label{sec:pareto}

A single utility presumes that the weights in \eqref{eq:utility} have been agreed. When they have not, let each stakeholder $s$ in a finite set $\mathcal S$ have a utility $U^{(s)}$, say that $q$ \emph{dominates} $p$ if $U^{(s)}(q)\ge U^{(s)}(p)$ for every $s$ with strict inequality for some $s$, and let each test return one noisy observation per stakeholder. In place of a favourite the team now keeps an \emph{archive} $A_k$: the tested prototypes that no other tested prototype dominates according to current beliefs.

\begin{proposition}
\label{prop:pareto}
Let Assumptions~\ref{as:finite}, \ref{as:local}, \ref{as:test} and~\ref{as:retention} hold for every stakeholder, with no ties in any $U^{(s)}$. Suppose that whenever a member of the archive has a dominating variant that has not been built, round $k$ builds such a variant with conditional probability at least $\delta_k$, where the $\delta_k$ are as in Assumption~\ref{as:gen}. Then almost surely the archive is eventually constant, it is the set of non-dominated members of $P_\infty$, and none of its members has a dominating variant.
\end{proposition}

\begin{proof}
As before $\Pk=P_\infty$ from some round on and all beliefs converge to the true utilities. With no ties, dominance according to beliefs eventually agrees with true dominance on the finite set $P_\infty$, so the archive settles on the non-dominated members of $P_\infty$. If a member $p$ had a dominating variant $q$, then $q\notin P_\infty$, because otherwise $p$ would not be in the archive. The generativity condition would then apply in every late round and a dominating variant would be built, contradicting $\Pk=P_\infty$.
\end{proof}

With one stakeholder the archive is the favourite and the proposition is Theorem~\ref{thm:converge}. With several, the loop does not resolve the disagreement. It delivers a set of mutually incomparable prototypes, each of which is locally unimprovable for everyone at once, and leaves the choice among them to negotiation. Archive-based evolutionary multi-objective methods behave in the same way \citep{laumanns2002}.

\subsection{How fast?}
\label{sec:rates}

Almost-sure convergence says nothing about time. The time has two components, the discovery of better prototypes and their recognition.

Discovery is governed by $\delta$. Suppose for a moment that tests are exact, so that the favourite is always the best prototype built so far, and that $\delta_k\ge\delta>0$. In each round in which the favourite is not yet locally optimal it improves with probability at least $\delta$, so the expected number of rounds until it is locally optimal is at most the expected number of improvements made on the way divided by $\delta$, a waiting-time bound in the spirit of the drift analysis of evolutionary algorithms \citep{heyao2001}. In the simulations below the best prototype tested improves about seven times in the first 80 rounds and fewer than three times after round 20.

Recognition is governed by the noise. With noisy tests the team must also come to believe that a better prototype is better. Telling apart two prototypes whose utilities differ by $\Delta$ requires a number of tests of the order of $\sigma^2/\Delta^2$ on each \citep{evendar2006, mannor2004}. Near the top of a landscape the differences are small. With $\sigma=0.1$ and prototypes that differ by one hundredth of a utility unit, a single comparison needs hundreds of tests, and the last comparisons in a project are for this reason the most expensive ones. The efficient response, known from ranking and selection, is to concentrate tests on the leading alternatives \citep{chen2000}.

\subsection{Stopping and reframing}
\label{sec:reframe}

Because the limit is local, more iteration cannot repair a poor frame. A team therefore faces two decisions that the theorem does not address: when to stop, and when to reframe. Both can be put on one scale. The value of another round within the frame is the \emph{expected improvement}
\begin{equation}
\label{eq:ei}
\mathrm{EI}_k=\max_p\E\big[(U(p)-X_k)^+\mid\filt_k\big],
\end{equation}
the criterion of Bayesian optimization \citep{jones1998} and a close relative of the knowledge gradient \citep{frazier2008}. The maximum is taken over the prototypes already tested and the ideas not yet built, an idea being judged by its rough estimate and by what the team has learned about prototypes in this frame. Because one noisy test does not reveal $U(p)$, the expected improvement is an upper bound on what a round can add. If a round costs $c_{\mathrm{iter}}$ in the same units, iterating is worthwhile at most while $\mathrm{EI}_k>c_{\mathrm{iter}}$. Simon's aspiration level then becomes a price: cheap tests justify long iteration, expensive ones an early stop \citep{simon1956}.

A reframe costs $c_{\mathrm{ref}}$ in the same units, the effort of going back to users and rebuilding context, and offers an uncertain gain $G_{\mathrm{ref}}$ over what the team already holds. A one-step comparison suggests reframing when
\begin{equation}
\label{eq:reframe}
\mathrm{EI}_k-c_{\mathrm{iter}}<G_{\mathrm{ref}}-c_{\mathrm{ref}} .
\end{equation}
The left side is computed from the team's beliefs. The right side contains the one quantity that cannot be derived from inside the current frame. It can be estimated as an expected improvement one level up. If what the team comes to hold in a frame is regarded as a draw from a normal distribution with mean $m$ and standard deviation $s$, learned from earlier projects and updated after each epoch, then
\begin{equation}
\label{eq:gref}
G_{\mathrm{ref}}=(m-X_k)\,\Phi(z)+s\,\phi(z),\qquad z=\frac{m-X_k}{s},
\end{equation}
with $\Phi$ and $\phi$ the standard normal distribution and density functions. Frames are then the arms of a bandit, iteration is exploitation and reframing is exploration in the sense of \citet{march1991}, and \eqref{eq:reframe} is the greedy policy for that bandit. It is a heuristic, and a lopsided one, since it sets the gain of a single round against the gain of a whole epoch in a new frame. Optimal search over alternatives with known reward distributions has the reservation-price structure of \citet{weitzman1979}, and the choice between trial-and-error learning and the parallel pursuit of several approaches is treated by \citet{sommerloch2004}. Section~\ref{sec:reframesim} compares the heuristic with simpler rules.

\section{Simulation study}
\label{sec:sim}

The simulations measure what the theorem leaves open: how far a team gets in a given number of rounds, which of the assumptions carry most weight at that horizon, and how the loop compares with other ways of spending the same tests. All code, seeds and raw outputs are in the repository listed at the end of the paper.

\subsection{Setup}
\label{sec:setup}

A prototype is a subset of $|F|=14$ features, so there are $2^{14}=16{,}384$ prototypes. This is small enough to compute the utility of every one, and therefore to know the optimum $U^*$ and the complete set of local optima exactly. We use three classes of landscape. The \emph{main landscape} is an instance of \eqref{eq:utility} with weights $(w_{\mathrm d},w_{\mathrm f},w_{\mathrm v})=(0.45,0.25,0.30)$: desirability rewards the coverage of latent user needs with diminishing returns, feasibility falls with build cost and with the cost of integrating pairs of features, and viability rewards business value and penalises scope. Its optimum is $U^*=0.763$, its mean utility $0.670$ and the standard deviation of utility across prototypes $0.046$. It has four local optima when variants differ in one feature and two when they may differ in two. The \emph{DVF family}, named after the three criteria, consists of 30 landscapes of the same form with independently drawn needs, values, costs and interactions (one to six local optima for one-feature variants, median three). The \emph{NK family} consists of 20 landscapes from Kauffman's NK model \citep{kauffman1993}, in which each feature's contribution depends on a number of other features, here 1, 2, 4 or 6, with five landscapes for each. These are far more rugged, with between 1 and 185 local optima for one-feature variants (median 38), and they share nothing with the main landscape beyond the size of the space. A result that holds on all three classes does not depend on how we chose to write down desirability, feasibility and viability.

A run starts from five random ideas. In each round the team picks at random one of its three \emph{leading prototypes}, those with the highest beliefs, and proposes five variants of it, each differing from it in one or two features. Among all the ideas it has not yet built, it builds the one that looks best on a rough estimate and tests it once. It then re-tests its three leading prototypes, and finally re-tests one further prototype. For this last test all prototypes are taken in rotation, so that each is tested again and again, as Assumption~\ref{as:test} requires, though at intervals that lengthen as their number grows; after 80 rounds about a third of the prototypes have still been tested only once. A round uses five tests, apart from the first two rounds, which use fewer. Tests follow \eqref{eq:test} with $\sigma=0.10$ and no bias. Beliefs are conjugate normal posterior means under an uninformative prior (mean $0.5$, variance $0.08$), confined to $[0,1]$. The rough estimate used for pre-screening is the true utility plus noise of standard deviation $0.10$, drawn once per idea. It is as informative as one baseline test and costs nothing, an advantage that Sections~\ref{sec:ablation} and~\ref{sec:baselines} remove.

Two choices in this setup shape everything that follows. The noise of a test is about twice the spread of utilities across the landscape, so a single test cannot tell an average prototype from a good one. We regard this as a fair picture of an early test with a handful of users, and Section~\ref{sec:ablation} halves and doubles it. A round, with one new prototype and four re-tests, is a substantial unit of work. A project of a few weeks may run five or ten such rounds, so we report results from 5 rounds upward and treat 80 rounds, about 400 tests, as a long project.

Utilities are reported as a percentage of the landscape's optimum $U^*$, and a difference of one \emph{point} is one percent of $U^*$. Intervals are 95\% confidence intervals; on the two families they are computed across landscapes from landscape means.

\subsection{The first 80 rounds}
\label{sec:dynamics}

Figure~\ref{fig:dynamics} shows the baseline on the main landscape. Progress is fast at first and slow afterwards. The shipped utility is 90.8\% of the optimum after 5 rounds, 92.1\% after 10, 93.3\% after 20 and 94.8\% after 80. At every one of these horizons the team has already tested something better than what it would ship: the best prototype found is worth 92.7\%, 94.2\%, 95.5\% and 97.5\%. After 80 rounds the team therefore loses 2.5 points to prototypes it has not found and 2.7 points to its own choice among those it has.

\begin{figure}[t]
\centering
\includegraphics[width=\columnwidth]{FigA_v5.png}
\caption{The baseline loop on the main landscape (250 runs; bands are 95\% confidence intervals). The team's belief about its favourite exceeds the optimum from round 18 on, while the favourite is worth less than the best prototype the team has tested. A prior calibrated to the landscape (dashed) removes most of the overestimate.}
\label{fig:dynamics}
\end{figure}

The balance between the two losses depends on the landscape. On the DVF family the shortfall in discovery after 80 rounds is 1.4 points and the loss in choosing 2.9. On the rugged NK family the order is reversed, at 3.6 and 2.7 points, because good prototypes are harder to find there. The cost of choosing is thus comparable to the cost of not finding on every class, and on the smoother landscapes it is the larger of the two.

The team's belief about its favourite behaves quite differently from the favourite's utility. It passes 100\% of the optimum at round 18 and ends at 102.4\%: the team believes it holds something better than the best prototype that exists. The overestimate at round 80 is $0.058\pm0.004$ in utility units, or 7.6\% of the optimum, and it is of similar size on the DVF family (7.3\%) and smaller on the NK family (3.8\%). This is the optimizer's curse \citep{smithwinkler2006}. The favourite is chosen for having the highest estimate, a high estimate is partly luck, and an uninformative prior does nothing to discount the luck. Replacing the prior by one matched to the main landscape (mean $0.670$, standard deviation $0.046$) cuts the overestimate to $0.014\pm0.002$ and leaves the shipped utility where it was (94.6\% against 94.8\%). Calibration corrects what the team believes about its result without improving the result.

\subsection{How far is the limit?}
\label{sec:horizon}

To see where the process is heading we ran it for up to 2{,}560 rounds and checked at each horizon whether the favourite is a local optimum, by looking it up in the enumerated set (Figure~\ref{fig:horizon}). On the main landscape it is one in 6 of 120 runs after 80 rounds and in 16 of 120 after 2{,}560 rounds, which is about 12{,}800 tests in a space of 16{,}384 prototypes. This landscape has only two local optima for two-feature variants, so reaching one nearly always means reaching the global optimum; with one-feature variants and four local optima the share after 2{,}560 rounds is 19\%. On the DVF family the share is 2\% after 80 rounds and 15\% after 1{,}280. On the NK family it is 26\% after 80 rounds and 72\% after 1{,}280, because four of the ten NK landscapes in this experiment have about twenty or more local optima even for two-feature variants, so that a local optimum is never far away.

\begin{figure*}[t]
\centering
\includegraphics[width=0.86\textwidth]{FigB_v5.png}
\caption{Short and long runs (main landscape: 120 runs; DVF and NK families: 10 landscapes each, 24 runs per landscape). \textbf{(a)} Utility shipped (solid) and best utility among tested prototypes (dashed). \textbf{(b)} Share of runs in which the favourite (solid) or the best tested prototype (dashed) is a local optimum, determined exactly by enumeration.}
\label{fig:horizon}
\end{figure*}

A team that stops after a few dozen rounds has in most cases not converged, and the value of its result is set by how quickly the loop climbs. Finishing below the global optimum is therefore not evidence of having reached a local one: on the main landscape nearly all runs finish below the global optimum and very few at a local optimum.

The obstacle is the ratio of test noise to the differences between prototypes, far more than the number of peaks. The dashed curves in Figure~\ref{fig:horizon} show this. After 2{,}560 rounds on the main landscape the best prototype tested is a local optimum in 88\% of runs and is worth 99.9\% of the optimum, while the favourite is worth 95.8\%: in these runs the gap between what has been found and what would be shipped has grown from 2.5 points after 80 rounds to 4.1. The reason is the quadratic cost of comparison noted in Section~\ref{sec:rates}. Late in a run the team holds thousands of prototypes within a few percent of each other, most of them tested once or twice, and the highest of several thousand noisy estimates rarely belongs to the best prototype. Where prototypes differ more, as on the NK family, the gap closes with time, from 2.8 points after 80 rounds to 0.6 after 1{,}280.

One remedy is to discount thin evidence by choosing on a lower confidence bound in place of the mean. A team that ships the prototype with the highest belief minus one posterior standard deviation gains little after 80 rounds (95.4\% against 95.0\% on the main landscape) and 2.4 points after 2{,}560 (98.2\% against 95.8\%). On the DVF family the gain after 1{,}280 rounds is 1.2 points, and on the NK family, where little is left to gain, 0.2.

\subsection{New prototypes or re-tests?}
\label{sec:allocation}

If choosing among prototypes is costly, the division of tests between new prototypes and old ones should make a difference. We gave every team exactly five tests in each of 60 rounds and varied how many of the five went to new prototypes, from one every other round to all five. The remaining tests were re-tests spread over the three leading prototypes. Ideas were variants of the leading prototypes in every arm, so that the arms differ in the use of tests and in nothing else (Figure~\ref{fig:allocation}; 120 runs on the main landscape and 16 on each landscape of the families, repeated with 5 and with 15 ideas per round).

\begin{figure*}[t]
\centering
\includegraphics[width=0.9\textwidth]{FigC_v5.png}
\caption{Shipped utility after 60 rounds of exactly five tests each, as a function of how many of the five tests go to new prototypes, for four levels of test noise (five ideas per round; error bars are 95\% confidence intervals; circles mark the best allocation).}
\label{fig:allocation}
\end{figure*}

The curves rise and then fall. Giving all five tests to new prototypes, so that nothing is ever tested twice, is the worst allocation in 23 of the 24 combinations of landscape class, noise level and number of ideas. Against the best allocation it costs 1.6 to 2.9 points on the main landscape, 1.4 to 6.5 on the DVF family and 1.3 to 11.5 on the NK family, and the cost grows with the noise. At the other end, building a new prototype only every second round is never the best allocation. It falls short of the best by up to one point on the main landscape and the DVF family and up to two on the NK family when tests are precise, and by at most 0.7 points at the baseline noise and above. The best allocation is one new prototype per round in 15 of the 24 combinations, two in five and three in four. It moves towards re-testing as tests get noisier: with precise tests ($\sigma=0.025$) two or three new prototypes per round are best, and at the baseline noise and above one is best on every class. Tripling the number of ideas per round changes none of this.

A team should therefore build something new in every round and, unless its tests are very precise, spend most of its tests on prototypes it has tested before. \citet{dow2010} found that building several prototypes before seeking feedback beats refining one at a time. Here too the gain comes from having alternatives to compare; what the simulations add is that the comparison is unreliable until the alternatives have been tested more than once.

\subsection{Relaxing the assumptions}
\label{sec:ablation}

Table~\ref{tab:ablation} changes one element of the baseline at a time, grouped by the assumption each change bears on.

\begin{table*}[t]
\caption{One change at a time: utility shipped after 80 rounds, as a percentage of $U^*$, with 95\% confidence half-widths (main landscape: 200 runs; families: 16 runs on each landscape). The baseline uses variants within two features, five ideas per round, full retention, one new prototype, three re-tests of leading prototypes and one rotating re-test per round, and $\sigma=0.10$.}
\label{tab:ablation}
\centering\small
\begin{tabular}{@{}llccc@{}}
\toprule
Assumption & Changed element & Main & DVF family & NK family\\
\midrule
Baseline & --- & $94.9\pm0.3$ & $95.7\pm0.5$ & $93.7\pm0.9$\\
\addlinespace[2pt]
Local ideation & variants within 1 feature & $94.5\pm0.3$ & $95.2\pm0.4$ & $92.5\pm1.1$\\
 & variants within 3 features & $94.7\pm0.3$ & $95.9\pm0.4$ & $93.3\pm0.8$\\
 & unbounded variants & $94.2\pm0.3$ & $94.9\pm0.5$ & $92.4\pm0.8$\\
\addlinespace[2pt]
Generativity & 1 idea per round & $94.5\pm0.3$ & $94.4\pm0.5$ & $89.6\pm1.2$\\
 & 2 ideas per round & $94.8\pm0.3$ & $95.2\pm0.4$ & $92.0\pm0.9$\\
 & 15 ideas per round & $94.8\pm0.3$ & $95.8\pm0.4$ & $94.0\pm0.9$\\
 & ideas built without pre-screening & $94.0\pm0.3$ & $94.3\pm0.6$ & $89.0\pm1.4$\\
\addlinespace[2pt]
Retention & only the favourite kept & $93.7\pm0.3$ & $93.8\pm0.6$ & $90.1\pm1.0$\\
 & best 3 kept & $94.9\pm0.2$ & $95.7\pm0.4$ & $94.0\pm0.9$\\
 & best 10 kept & $94.9\pm0.3$ & $96.0\pm0.4$ & $93.4\pm0.9$\\
\addlinespace[2pt]
Testing & all five tests on new prototypes & $92.6\pm0.4$ & $92.0\pm0.7$ & $85.4\pm1.2$\\
 & no rotating re-test & $94.7\pm0.3$ & $95.7\pm0.4$ & $93.0\pm0.9$\\
 & noise $\sigma=0.05$ & $95.9\pm0.2$ & $97.5\pm0.4$ & $95.4\pm1.0$\\
 & noise $\sigma=0.20$ & $93.5\pm0.3$ & $93.6\pm0.5$ & $88.6\pm1.3$\\
\bottomrule
\end{tabular}
\end{table*}

The reach of ideation has little effect within limits. Narrowing it to one-feature variants costs 0.4 points on the main landscape and the DVF family and 1.2 on the NK family, and widening it to three features changes nothing measurable. Removing the bound altogether, so that a variant may differ from its parent in any number of features, lowers the result by 0.7 to 1.3 points. In the limit unbounded ideation reaches the global optimum, but within 80 rounds it spends tests on distant prototypes that are mostly poor.

The supply of ideas is important on rugged landscapes and saturates quickly elsewhere. A team that produces one idea per round ships 0.4 points less than the baseline on the main landscape and 4.0 points less on the NK family, and fifteen ideas per round are no better than five on any class. Which idea is built counts for more than how many there are to choose from. Building one of the current round's ideas at random, in place of the idea that looks best among all those not yet built, costs 0.9, 1.4 and 4.7 points on the three classes.

Retention has a specific meaning in practice. A team that remembers only its current favourite loses 1.2 to 3.6 points, and remembering three prototypes recovers all of it. What the favourite needs is rivals to be compared with; a complete archive is not required at this horizon.

Testing has the largest effects in the table. Spending all five tests of a round on new prototypes costs 2.3, 3.7 and 8.2 points, although such a team finds prototypes as good as the baseline team does; it simply cannot tell which they are. Dropping only the rotating re-test of older prototypes, which the theorem needs for beliefs to become correct, costs at most 0.6 points within 80 rounds. Halving the noise of a test is worth 1.0 to 1.8 points and doubling it costs 1.4 to 5.1.

\begin{figure*}[t]
\centering
\includegraphics[width=0.9\textwidth]{FigD_v5.png}
\caption{Change in shipped utility after 80 rounds when each round's re-test of the favourite is inflated by $b$, at three levels of test noise (main landscape: 240 runs per point; families: 16 runs per landscape; error bars are 95\% confidence intervals). The dashed line marks $b=\sigma$.}
\label{fig:bias}
\end{figure*}

Figure~\ref{fig:bias} adds a confirmation bias: in each round the re-test of the current favourite is inflated by a constant $b$. At the baseline noise of $0.10$ no loss is detectable on any class up to $b=0.10$; at $b=0.20$ the loss is 1.0 to 3.5 points and at $b=0.30$ it is 3.7 to 8.0. The point at which losses begin moves with the noise. With more precise tests ($\sigma=0.05$) losses are visible from $b=0.075$ and reach 1.1 to 2.5 points at $b=0.10$, and with noisier tests ($\sigma=0.20$) nothing is lost up to $b=0.20$. A bias protects the favourite against challengers, and it changes the outcome only when it is large enough to reverse comparisons that the evidence would otherwise have decided. In these experiments that means a bias about as large as the standard deviation of a test, which is many times larger than the differences among the leading prototypes. A team with precise tests is the more exposed to its own bias, because less noise is left in which a challenger can be lucky. Proposition~\ref{prop:bias} gives a guarantee that weakens in proportion to the bias, whereas the typical loss has a threshold.

\subsection{Comparison with other policies}
\label{sec:baselines}

To see whether the loop is a good way to spend a testing budget, we compared it with five alternatives given 300 tests, the same variants, the same test noise and the same beliefs (Table~\ref{tab:baselines}): random search among variants; hill-climbing that tests each one-feature variant of its current point several times and moves to the best; an $\varepsilon$-greedy policy that with probability $\varepsilon$ tests a new variant of the favourite and otherwise re-tests the favourite; an upper-confidence-bound (UCB) policy \citep{auer2002}; and Thompson sampling \citep{russo2018}, the last two adding a new variant of the favourite at regular intervals. The loop screens its ideas before building one and these policies do not. We therefore also ran the loop without pre-screening, building one of the current round's ideas at random, and the UCB policy with pre-screening, building the best of five fresh variants of the favourite on the same rough estimate. The parameters of each alternative were tuned on ten DVF and eight NK landscapes, and all policies were then scored on landscapes and seeds not used for tuning. The loop was run with its baseline settings for 60 rounds, which is 297 tests.

\begin{table*}[t]
\caption{Utility shipped after 300 tests, as a percentage of $U^*$, on landscapes not used for tuning (main landscape: 250 runs; families: 20 landscapes each, 30 runs per landscape; $\pm$ is the 95\% confidence half-width). ``Diff.'' is the mean paired difference from the loop across landscapes.}
\label{tab:baselines}
\centering\small
\begin{tabular}{@{}lccccc@{}}
\toprule
& Main & \multicolumn{2}{c}{DVF family} & \multicolumn{2}{c}{NK family}\\
\cmidrule(lr){3-4}\cmidrule(lr){5-6}
Policy & shipped & shipped & diff. & shipped & diff.\\
\midrule
Random search & $91.0\pm0.4$ & $87.9\pm0.8$ & $-7.4\pm0.8$ & $80.2\pm0.9$ & $-12.0\pm0.9$\\
Hill-climbing & $91.0\pm0.5$ & $88.7\pm0.8$ & $-6.6\pm0.8$ & $83.8\pm1.4$ & $-8.4\pm1.6$\\
$\varepsilon$-greedy & $94.0\pm0.3$ & $94.3\pm0.6$ & $-1.1\pm0.3$ & $89.9\pm1.1$ & $-2.3\pm0.6$\\
UCB & $94.5\pm0.3$ & $94.7\pm0.5$ & $-0.6\pm0.3$ & $90.2\pm1.1$ & $-2.1\pm0.5$\\
Thompson sampling & $93.9\pm0.3$ & $94.4\pm0.7$ & $-0.9\pm0.4$ & $89.6\pm1.1$ & $-2.7\pm0.6$\\
\addlinespace[2pt]
UCB with pre-screening & $95.0\pm0.2$ & $96.1\pm0.4$ & $+0.8\pm0.2$ & $93.6\pm1.0$ & $+1.4\pm0.5$\\
Loop without pre-screening & $93.8\pm0.3$ & $93.9\pm0.6$ & $-1.5\pm0.5$ & $87.5\pm1.2$ & $-4.7\pm0.7$\\
Loop & $94.7\pm0.3$ & $95.3\pm0.4$ & --- & $92.2\pm0.9$ & ---\\
\bottomrule
\end{tabular}
\end{table*}

The table separates the policy for spending tests from the pre-screening of ideas. As to the policy, random search and hill-climbing are 4 to 12 points behind the loop, because a policy that does not concentrate its tests on promising prototypes wastes most of them. Among the policies that do concentrate their tests the differences are much smaller, and the loop is not the best of them. Without pre-screening it is level with $\varepsilon$-greedy and Thompson sampling and 0.7 points behind UCB on the main landscape, 0.4 to 0.9 points behind the three bandit policies on the DVF family and 2.1 to 2.7 points behind them on the NK family. With pre-screening on both sides, UCB leads the loop by 0.3 points on the main landscape, which is within noise, and by 0.8 and 1.4 points on the families, where it wins on nine landscapes in ten.

Pre-screening is worth 0.9, 1.5 and 4.7 points to the loop and 0.5, 1.4 and 3.5 points to UCB, which on every class exceeds the difference between the two policies. A fixed routine of one new prototype and four re-tests per round is thus a good way to spend a testing budget. With pre-screening it is within a point and a half of a bandit policy tuned for the purpose, it is far ahead of search that does not re-test, and the pre-screening of ideas contributes more than the choice of routine.

\subsection{Reframing}
\label{sec:reframesim}

The last experiment concerns the outer level. A world contains six frames, each a different DVF landscape drawn from a pool of 40. A team has 150 rounds. It works in one frame at a time with the loop of Section~\ref{sec:setup}, here without the rotating re-test, and a reframe costs $c_{\mathrm{ref}}$ rounds and moves it to a frame it has not visited, where it starts again. When the rounds are spent, the team ships whichever of the favourites of the frames it has visited it believes to be best, and the payoff is that prototype's true utility. We compare a team that never reframes, teams that leave each frame after a fixed number of rounds, and teams that follow criterion~\eqref{eq:reframe} after a minimum stay of ten rounds. No policy reframes unless at least ten rounds would remain afterwards. The criterion uses the team's beliefs only. A round costs $0.002$ in utility units and a reframe that many times its length in rounds, the expected improvement of an idea is computed from its rough estimate shrunk towards the mean of the team's beliefs in the frame, and $G_{\mathrm{ref}}$ follows \eqref{eq:gref} under three priors for what a frame yields: a default (mean $0.75$, standard deviation $0.05$), one calibrated to the beliefs teams actually hold after 75 rounds in these frames ($0.74$ and $0.04$), and an optimistic one ($0.79$ and $0.04$). The prior is widened by a frame-to-frame spread of $0.03$ and updated with the belief the team holds each time it leaves a frame. Every policy sees the same 240 worlds at each of four costs.

\begin{table}[t]
\caption{True utility of the prototype shipped after 150 rounds in a world of six frames, by reframing policy and by the cost of a reframe in rounds (240 paired runs; 95\% half-widths are at most $0.004$). The last column (refr.) is the mean number of reframes per run at a cost of 10 rounds; the criterion rows differ in the prior used for $G_{\mathrm{ref}}$. Bold marks the best policy at each cost.}
\label{tab:reframe}
\centering\small
\setlength{\tabcolsep}{3.4pt}
\begin{tabular}{@{}lccccc@{}}
\toprule
& \multicolumn{4}{c}{cost of a reframe (rounds)} & \\
\cmidrule(lr){2-5}
Policy & 4 & 10 & 20 & 35 & refr.\\
\midrule
Never & 0.694 & 0.694 & 0.694 & 0.694 & 0\\
\addlinespace[2pt]
Every 15 rounds & 0.700 & 0.697 & 0.693 & 0.694 & 5.0\\
Every 25 rounds & \textbf{0.703} & 0.698 & 0.698 & 0.698 & 4.0\\
Every 40 rounds & 0.702 & \textbf{0.704} & \textbf{0.703} & 0.698 & 2.0\\
Once, after 75 rounds & 0.701 & 0.702 & 0.699 & 0.699 & 1.0\\
\addlinespace[2pt]
Criterion, default & 0.703 & 0.701 & 0.694 & 0.699 & 1.2\\
Criterion, calibrated & 0.702 & 0.702 & 0.697 & \textbf{0.699} & 1.2\\
Criterion, optimistic & 0.703 & 0.701 & 0.693 & 0.695 & 1.9\\
\bottomrule
\end{tabular}
\end{table}

A single reframe at the midpoint is better than none at every cost, clearly so when a reframe is cheap (Table~\ref{tab:reframe}). The paired gain is $0.007\pm0.004$ and $0.008\pm0.004$ in utility at costs of 4 and 10 rounds and $0.004\pm0.004$ and $0.005\pm0.004$ at costs of 20 and 35, about one percent of what the team ships. Reframing more often adds little. Leaving a frame every 40 rounds is better than the single reframe by $0.001$ to $0.005$ at costs up to 20 rounds, mostly within the confidence intervals, and leaving every 15 rounds is no better than never reframing once a reframe costs 10 rounds or more.

The criterion does about as well as the single scheduled reframe and no better. With the calibrated prior it differs from it by less than the confidence interval at all four costs, and with the default prior at three of them; at a cost of 20 rounds the default prior is worse by $0.004\pm0.004$. In most runs it fires at the earliest permitted moment (with the default prior and a cost of 10 rounds the median is round 10, and nine in ten reframes occur by round 40), so it works as an early exit from frames that look poor against the prior and seldom as a signal that iteration has run dry. The expected improvement of a further round does not decline along an epoch as one might expect. New, thinly tested rivals keep it up: it averages $0.011$ in rounds 10 to 20 and $0.011$ in rounds 40 to 75, and it exceeds the cost of a round in 86\% of all rounds, so on its own it would seldom tell the team to stop.

In this model most of the value of reframing lies in doing it once, and a rule as crude as going back to the problem halfway through captures that value. Criterion~\eqref{eq:reframe} states what the decision depends on, but as an operating rule it adds nothing to a fixed plan.

\section{Discussion}
\label{sec:discussion}

We first answer the question in the title, then consider what follows for tools and where the model stops.

\subsection{When design thinking works}

In the limit, the iterate--test loop works when the five assumptions hold, and each of them names a practice: keep the brief fixed while iterating, stay within reach of what has been built, keep producing variants, test in ways that can contradict the team and test favourites again, and keep what has been tried. When these hold, improvement follows from the structure of the process and does not depend on luck or talent. This is consistent with Liedtka's proposal that the method works by disciplining judgement \citep{liedtka2015}. Table~\ref{tab:conditions} sets the assumptions beside what each asks of a team and what a departure from it cost in the simulations.

\begin{table*}[t]
\caption{The five assumptions of Theorem~\ref{thm:converge}: what each means in practice, what the theory says when it fails, and what a departure in its direction cost after 80 rounds (range over the three landscape classes, in points of $U^*$; from Table~\ref{tab:ablation} and Figure~\ref{fig:bias}).}
\label{tab:conditions}
\centering\small
\begin{tabularx}{\textwidth}{@{}p{2.6cm}XXp{4.2cm}@{}}
\toprule
Assumption & In practice & When it fails & Simulated departure and its cost\\
\midrule
Finite frame & The brief and the vocabulary of features stay fixed during an epoch & No fixed space in which an optimum is defined & Not simulated\\
Local ideation & New ideas are variants of what has been built & With unbounded jumps the limit is global but progress is slower & Unbounded variants: $-0.7$ to $-1.3$\\
Persistent generativity & The team keeps proposing variants, even late & The process can stop short of a local optimum & One idea per round in place of five: $-0.4$ to $-4.0$\\
Unbiased, repeated testing & Tests can contradict the team; favourites are tested again & A lucky prototype can stay on top; under bias $B$ the result is only $2B$-locally optimal & No prototype tested twice: $-2.3$ to $-8.2$. Bias $b=0.2$: $-1.0$ to $-3.5$\\
Retention & Rivals to the favourite are kept & The best prototype found can be lost & Only the favourite kept: $-1.2$ to $-3.6$\\
\bottomrule
\end{tabularx}
\end{table*}

At the horizon of an actual project the answer shifts. Design thinking is usually presented as a discipline of generating: many ideas, many cheap prototypes, an early move from talking to building \citep{brown2008, dschool2010}. In our simulations generation is only half of the difficulty. Teams find good prototypes early, and at every horizon they have tested something better than what they would ship. Their confidence in what they would ship is inflated, and a prototype tested once is a poor basis for a decision. The pattern that an emphasis on building invites, many different prototypes each shown to users once, is the worst allocation we found. Experiments on idea selection point the same way: groups that generate more ideas do not select better ones, and their selections are little better than chance \citep{rietzschel2006}, while the quality of the best idea a group ends with depends among other things on its ability to discern quality \citep{girotra2010}.

Four recommendations follow. A team should screen its ideas before building, since a rough estimate of which idea is most promising was worth more in our comparisons than the difference between the loop's fixed routine and a tuned bandit policy. It should treat its estimate of its favourite concept as an overestimate; the optimizer's curse grows with the number of alternatives chosen among \citep{smithwinkler2006}. It should build something new in every round and give most of its remaining tests to its two or three leading concepts, preferring a well-tested concept to one that looked excellent once. And it should worry about confirmation bias most when its tests are precise, since a bias smaller than the noise of a test changed nothing in our simulations and a bias well above the noise cost several points.

On reframing the simulations support a plain rule. Iteration cannot escape a poor frame, a team that reframes once does better than one that never does, and a rule that watches the expected improvement of further iteration added nothing to a fixed plan. A team should plan for at least one return to the problem definition and should not expect iteration to signal when it is due.

\subsection{Implications for tools}

Shared whiteboards and similar tools help design teams generate and arrange ideas. The model suggests that evaluation is the scarcer resource, and it specifies what a tool that supported evaluation would do: keep a belief and an evidence count for every concept tested, shrink estimates that rest on little evidence, show the team how much of its favourite's lead is likely to be selection, and propose the next test by its expected improvement. The same reasoning applies to generative AI as a source of ideas. In a story-writing experiment, access to ideas from a language model raised the rated creativity of individual stories and made the stories more alike \citep{doshi2024}. In the model, more ideas per round help only on rugged landscapes, and ideas that resemble each other amount to a smaller neighbourhood. A language model would be more useful in pre-screening, if it can rank ideas even roughly before any is built.

\subsection{Limitations}

The model is normative. It describes a team that keeps beliefs and revises them consistently, which real teams do informally at best, and its predictions concern what happens as a team's practice comes closer to the assumptions or moves away from them.

The frame is taken as given. The model is silent on how a point of view is formed and on what makes one frame better than another, which is where much of the skill in design lies and where problems are wicked in the sense of \citet{rittel1973} and \citet{buchanan1992}, without a definitive formulation or a stopping rule. Holding the frame fixed during an epoch also sits uneasily with the view that problem and solution evolve together; Corollary~\ref{cor:drift} accommodates a problem that shifts within bounds and nothing more. The reframing experiment treats frames as interchangeable draws from a population.

Within the frame, the simplification with the largest consequences concerns what a test returns. In the model a test returns a number. In practice it also returns reasons: a user test shows why a prototype fails, and the reasons direct the next idea. The model therefore understates what a single test of a new prototype is worth, and the allocation result should be read with that in mind. Independent beliefs work in the same direction: a test of one prototype teaches something about its relatives, and correlated beliefs would speed up recognition and call for the correlated form of the knowledge gradient \citep{frazier2009}. The representation is restrictive as well. A prototype is a set of features drawn from a fixed list, which leaves out changes of architecture or of a continuous parameter, and the discovery of a new kind of feature in mid-epoch.

The simulations use synthetic landscapes with fourteen features. The qualitative results hold across three classes of landscape of very different ruggedness, but we have not varied the size of the space, and the magnitudes should not be read as estimates for real projects. The threshold for bias was located at three noise levels only, the reframing experiment uses one class of landscape and one cost per round, and the rough pre-screening of ideas, which is free in our simulations, would cost something in practice.

\section{Testing the model on real teams}
\label{sec:empirical}

The results so far are statements about a model. This section describes what would have to be observed in real projects to test it.

\subsection{Predictions}

In the vocabulary of practice a prototype is a concept, and we use that word here. That teams sometimes carry forward a concept that is not their best is predicted by any account in which evaluation is noisy, and observing it would not distinguish our model from others. Three more specific consequences would. Re-testing should reduce the cost of choosing: the shortfall of the concept a team carries forward, relative to the best concept it tested, should be larger in teams that tested each concept once than in teams that returned to their leading concepts, and the team's own assessment of its chosen concept should exceed an independent assessment by more when it chose among more candidates. Allocation should have an interior optimum: among teams with similar testing effort, outcomes should be best at an intermediate share of tests spent on new concepts and should fall off sharply as that share approaches one. Bias should have a threshold: leading questions and tests confined to the favourite should be associated with worse outcomes only where the bias they introduce is large relative to the variability of test results, and so more often in teams whose tests are precise.

Each of these could fail. If teams turned out to carry forward their best-rated concept nearly always, the difficulty of choosing that dominates our simulations would be an artefact of the model's independent beliefs. If outcomes rose steadily with the number of distinct prototypes built, the allocation result would be wrong where it counts. The prediction about reframing, that teams which return to the problem definition at least once do better, cannot be tested on observational records at all, because teams reframe when they are in trouble; it requires an experiment in which a reframe is scheduled for some teams and not for others.

\subsection{Measurement}

The quantities involved can be obtained from recorded design sessions by protocol analysis, which has a long history in design research \citep{cross1996}. A session is segmented into moves \citep{goldschmidt2014}, and each move is labelled as proposing a concept, building one, testing one, or returning to the problem definition. For each test the coder records which concept it concerned and the direction and strength of the evidence on an ordinal scale, together with markers of likely bias such as a leading question or a test of the favourite alone. Two coders should label an overlapping sample and report a chance-corrected measure of agreement such as Krippendorff's $\alpha$ \citep{krippendorff2018}.

A study by \citet{sung2019} shows both that such coding supports a sequential analysis and where its limits lie. They coded nine think-aloud sessions, in which triads of fourth-grade pupils worked on a paper design task, into seven cognitive strategies (699 coded events, with 85\% raw agreement on four sessions coded twice) and analysed the 505 transitions between strategies by lag-sequential analysis \citep{bakeman2011}. Generating an idea (their code \emph{Designing}) was followed by anticipating its consequences (\emph{Predicting}) more often than chance would give, in 33 of the 147 transitions out of generating ($z=3.66$), and predicting was followed by a return to generating in 39 of 52 cases ($z=5.17$). Sketching an idea (\emph{Modeling}), by contrast, was followed directly by further generating less often than chance ($z=-4.47$). The cycle of generating and predicting that the study finds takes place in thought and is close to what we call pre-screening. The pupils built and tested nothing, the sessions were pooled, and the authors caution against generalising from them, so the study bears on the feasibility of the measurement and not on the loop of building and testing.

Three features of a confirmatory study are essential. The unit of analysis is the framing epoch, because the model's claims hold within a frame. The outcome must be measured independently of the team's own beliefs, by judges who rate every concept the team tested without knowing which one it chose; the consensual assessment technique gives a procedure and a reliability check for such ratings \citep{amabile1982}, and since the ratings are themselves noisy, several judges are needed for each concept. Comparisons should be made within epochs where possible. For a paired comparison with a standardised difference of 0.3, about one hundred epochs give a power of 0.9 at a one-sided level of 5\%.

\section{Conclusion}

Within a fixed framing of a problem, the loop of ideation, prototyping and testing converges to a locally optimal prototype when the frame is finite, ideation is local, improving ideas keep coming at a non-summable rate, tests are unbiased and repeated, and tested prototypes are retained. The guarantee weakens in proportion to a bounded bias or a bounded drift of the criteria, and with several stakeholders it becomes convergence to a locally Pareto-optimal set. The theorem identifies the practices on which the reliability of iteration rests.

The simulations show that the limit describes the direction of the process more than the position of a team that stops. Except on very rugged landscapes few runs reach a local optimum, and the result depends as much on how well the team can choose among what it has built as on what it has found. A prototype tested once is a poor basis for a decision, the team's confidence in its favourite is systematically too high, confirmation bias does harm once it exceeds the noise in the tests, and one reframe captures most of what reframing has to offer. A plain routine of one new prototype and four re-tests per round comes close to a bandit policy tuned for the task when both screen their ideas, and the pre-screening is worth more than the difference between them.

Design thinking works, in the sense made precise here, when a team keeps its options, generates near them, screens what it builds, and tests its leading prototypes more than once. Whether teams that do these things obtain better outcomes than teams that do not is a question for field data, and Section~\ref{sec:empirical} says which data would settle it.

\section*{Data and code availability}
The code that produces every figure, table and number in this paper, together with the raw outputs, is available at \url{https://github.com/pawelkuras/design-thinking-formalization} and will be deposited in a permanent archive on acceptance. The simulation study is produced by \texttt{engine.py}, \texttt{studies.py}, \texttt{reframing3.py} and \texttt{figures\_v5.py}; \texttt{bash run\_all.sh v5} runs them from scratch in the repository. The experiments use fixed random seeds and require only \texttt{numpy} and \texttt{matplotlib}.

\section*{CRediT authorship contribution statement}
\textbf{Pawe\l{} Kura\'s:} Conceptualization, Methodology, Formal analysis, Software, Writing -- original draft. \textbf{Adrian Michalski:} Formal analysis, Validation, Writing -- review \& editing. \textbf{Alicja Gerka:} Conceptualization, Investigation, Writing -- review \& editing. \textbf{Patryk Organi\'sciak:} Formal analysis, Validation, Writing -- review \& editing.

\section*{Declaration of competing interest}
The authors declare no competing financial interests or personal relationships that could have influenced the work reported in this paper.

\section*{Declaration of generative AI use}
A large-language-model assistant (Claude, Anthropic) was used extensively in preparing this manuscript: to draft and revise the text, to write and run the simulation code, to propose and carry out the simulation studies of Section~\ref{sec:sim}, to restate and check the proofs, to review successive drafts, and to verify the bibliography against publishers' records. The authors defined the research question and the model, directed and reviewed this work, and are responsible for the content of the paper.

\begin{thebibliography}{99}
\setlength{\itemsep}{0pt}\setlength{\parsep}{0pt}\small

\bibitem[Altman et~al.(2018)]{altman2018} Altman, M., Huang, T. T. K., \& Breland, J. Y. (2018). Design thinking in health care. \textit{Preventing Chronic Disease, 15}, E117. \doi{10.5888/pcd15.180128}

\bibitem[Amabile(1982)]{amabile1982} Amabile, T. M. (1982). Social psychology of creativity: A consensual assessment technique. \textit{Journal of Personality and Social Psychology, 43}(5), 997--1013. \doi{10.1037/0022-3514.43.5.997}

\bibitem[Audibert et~al.(2010)]{audibert2010} Audibert, J.-Y., Bubeck, S., \& Munos, R. (2010). Best arm identification in multi-armed bandits. In A. T. Kalai \& M. Mohri (Eds.), \textit{Proceedings of the 23rd Conference on Learning Theory (COLT 2010)} (pp. 41--53). Omnipress.

\bibitem[Auer et~al.(2002)]{auer2002} Auer, P., Cesa-Bianchi, N., \& Fischer, P. (2002). Finite-time analysis of the multiarmed bandit problem. \textit{Machine Learning, 47}(2--3), 235--256. \doi{10.1023/A:1013689704352}

\bibitem[Bakeman \& Quera(2011)]{bakeman2011} Bakeman, R., \& Quera, V. (2011). \textit{Sequential Analysis and Observational Methods for the Behavioral Sciences}. Cambridge University Press. \doi{10.1017/CBO9781139017343}

\bibitem[Brown(2008)]{brown2008} Brown, T. (2008). Design thinking. \textit{Harvard Business Review, 86}(6), 84--92.

\bibitem[Brown(2009)]{brown2009} Brown, T. (2009). \textit{Change by Design: How Design Thinking Transforms Organizations and Inspires Innovation}. HarperBusiness.

\bibitem[Buchanan(1992)]{buchanan1992} Buchanan, R. (1992). Wicked problems in design thinking. \textit{Design Issues, 8}(2), 5--21. \doi{10.2307/1511637}

\bibitem[Carlgren et~al.(2016)]{carlgren2016} Carlgren, L., Rauth, I., \& Elmquist, M. (2016). Framing design thinking: The concept in idea and enactment. \textit{Creativity and Innovation Management, 25}(1), 38--57. \doi{10.1111/caim.12153}

\bibitem[Chen et~al.(2000)]{chen2000} Chen, C.-H., Lin, J., Y{\"u}cesan, E., \& Chick, S. E. (2000). Simulation budget allocation for further enhancing the efficiency of ordinal optimization. \textit{Discrete Event Dynamic Systems, 10}(3), 251--270. \doi{10.1023/A:1008349927281}

\bibitem[Cross et~al.(1996)]{cross1996} Cross, N., Christiaans, H., \& Dorst, K. (Eds.). (1996). \textit{Analysing Design Activity}. Wiley.

\bibitem[Dahan \& Mendelson(2001)]{dahan2001} Dahan, E., \& Mendelson, H. (2001). An extreme-value model of concept testing. \textit{Management Science, 47}(1), 102--116. \doi{10.1287/mnsc.47.1.102.10666}

\bibitem[DeGroot(2004)]{degroot2004} DeGroot, M. H. (2004). \textit{Optimal Statistical Decisions}. Wiley-Interscience. (Original work published 1970)

\bibitem[Design Council(2007)]{designcouncil2007} Design Council. (2007). \textit{Eleven Lessons: Managing Design in Eleven Global Brands. A Study of the Design Process}. Design Council.

\bibitem[Dorst \& Cross(2001)]{dorstcross2001} Dorst, K., \& Cross, N. (2001). Creativity in the design process: Co-evolution of problem--solution. \textit{Design Studies, 22}(5), 425--437. \doi{10.1016/S0142-694X(01)00009-6}

\bibitem[Dorst(2011)]{dorst2011} Dorst, K. (2011). The core of `design thinking' and its application. \textit{Design Studies, 32}(6), 521--532. \doi{10.1016/j.destud.2011.07.006}

\bibitem[Doshi \& Hauser(2024)]{doshi2024} Doshi, A. R., \& Hauser, O. P. (2024). Generative AI enhances individual creativity but reduces the collective diversity of novel content. \textit{Science Advances, 10}(28), eadn5290. \doi{10.1126/sciadv.adn5290}

\bibitem[Dow et~al.(2010)]{dow2010} Dow, S. P., Glassco, A., Kass, J., Schwarz, M., Schwartz, D. L., \& Klemmer, S. R. (2010). Parallel prototyping leads to better design results, more divergence, and increased self-efficacy. \textit{ACM Transactions on Computer-Human Interaction, 17}(4), Article 18. \doi{10.1145/1879831.1879836}

\bibitem[Durrett(2019)]{durrett2019} Durrett, R. (2019). \textit{Probability: Theory and Examples} (5th ed.). Cambridge University Press. \doi{10.1017/9781108591034}

\bibitem[Erat \& Kavadias(2008)]{erat2008} Erat, S., \& Kavadias, S. (2008). Sequential testing of product designs: Implications for learning. \textit{Management Science, 54}(5), 956--968. \doi{10.1287/mnsc.1070.0784}

\bibitem[Even-Dar et~al.(2006)]{evendar2006} Even-Dar, E., Mannor, S., \& Mansour, Y. (2006). Action elimination and stopping conditions for the multi-armed bandit and reinforcement learning problems. \textit{Journal of Machine Learning Research, 7}, 1079--1105.

\bibitem[Frazier et~al.(2008)]{frazier2008} Frazier, P. I., Powell, W. B., \& Dayanik, S. (2008). A knowledge-gradient policy for sequential information collection. \textit{SIAM Journal on Control and Optimization, 47}(5), 2410--2439. \doi{10.1137/070693424}

\bibitem[Frazier et~al.(2009)]{frazier2009} Frazier, P., Powell, W., \& Dayanik, S. (2009). The knowledge-gradient policy for correlated normal beliefs. \textit{INFORMS Journal on Computing, 21}(4), 599--613. \doi{10.1287/ijoc.1080.0314}

\bibitem[Gero \& Kannengiesser(2004)]{gero2004} Gero, J. S., \& Kannengiesser, U. (2004). The situated function--behaviour--structure framework. \textit{Design Studies, 25}(4), 373--391. \doi{10.1016/j.destud.2003.10.010}

\bibitem[Girotra et~al.(2010)]{girotra2010} Girotra, K., Terwiesch, C., \& Ulrich, K. T. (2010). Idea generation and the quality of the best idea. \textit{Management Science, 56}(4), 591--605. \doi{10.1287/mnsc.1090.1144}

\bibitem[Goldschmidt(2014)]{goldschmidt2014} Goldschmidt, G. (2014). \textit{Linkography: Unfolding the Design Process}. MIT Press. \doi{10.7551/mitpress/9455.001.0001}

\bibitem[Gutjahr \& Pflug(1996)]{gutjahr1996} Gutjahr, W. J., \& Pflug, G. Ch. (1996). Simulated annealing for noisy cost functions. \textit{Journal of Global Optimization, 8}(1), 1--13. \doi{10.1007/BF00229298}

\bibitem[Hasso Plattner Institute of Design(2010)]{dschool2010} Hasso Plattner Institute of Design at Stanford. (2010). \textit{An Introduction to Design Thinking: Process Guide}. Stanford University, d.school.

\bibitem[Hatchuel \& Weil(2009)]{hatchuel2009} Hatchuel, A., \& Weil, B. (2009). C--K design theory: An advanced formulation. \textit{Research in Engineering Design, 19}(4), 181--192. \doi{10.1007/s00163-008-0043-4}

\bibitem[He \& Yao(2001)]{heyao2001} He, J., \& Yao, X. (2001). Drift analysis and average time complexity of evolutionary algorithms. \textit{Artificial Intelligence, 127}(1), 57--85. \doi{10.1016/S0004-3702(01)00058-3}

\bibitem[Homem-de-Mello(2003)]{homem2003} Homem-de-Mello, T. (2003). Variable-sample methods for stochastic optimization. \textit{ACM Transactions on Modeling and Computer Simulation, 13}(2), 108--133. \doi{10.1145/858481.858483}

\bibitem[Jansson \& Smith(1991)]{jansson1991} Jansson, D. G., \& Smith, S. M. (1991). Design fixation. \textit{Design Studies, 12}(1), 3--11. \doi{10.1016/0142-694X(91)90003-F}

\bibitem[Jin \& Branke(2005)]{jin2005} Jin, Y., \& Branke, J. (2005). Evolutionary optimization in uncertain environments---a survey. \textit{IEEE Transactions on Evolutionary Computation, 9}(3), 303--317. \doi{10.1109/TEVC.2005.846356}

\bibitem[Johansson-Sk{\"o}ldberg et~al.(2013)]{johansson2013} Johansson-Sk{\"o}ldberg, U., Woodilla, J., \& \c{C}etinkaya, M. (2013). Design thinking: Past, present and possible futures. \textit{Creativity and Innovation Management, 22}(2), 121--146. \doi{10.1111/caim.12023}

\bibitem[Jones et~al.(1998)]{jones1998} Jones, D. R., Schonlau, M., \& Welch, W. J. (1998). Efficient global optimization of expensive black-box functions. \textit{Journal of Global Optimization, 13}(4), 455--492. \doi{10.1023/A:1008306431147}

\bibitem[Kauffman \& Levin(1987)]{kauffmanlevin1987} Kauffman, S., \& Levin, S. (1987). Towards a general theory of adaptive walks on rugged landscapes. \textit{Journal of Theoretical Biology, 128}(1), 11--45. \doi{10.1016/S0022-5193(87)80029-2}

\bibitem[Kauffman \& Weinberger(1989)]{kauffmanweinberger1989} Kauffman, S. A., \& Weinberger, E. D. (1989). The NK model of rugged fitness landscapes and its application to maturation of the immune response. \textit{Journal of Theoretical Biology, 141}(2), 211--245. \doi{10.1016/S0022-5193(89)80019-0}

\bibitem[Kauffman(1993)]{kauffman1993} Kauffman, S. A. (1993). \textit{The Origins of Order: Self-Organization and Selection in Evolution}. Oxford University Press.

\bibitem[Keeney \& Raiffa(1976)]{keeney1976} Keeney, R. L., \& Raiffa, H. (1976). \textit{Decisions with Multiple Objectives: Preferences and Value Tradeoffs}. Wiley.

\bibitem[Knudsen \& Levinthal(2007)]{knudsen2007} Knudsen, T., \& Levinthal, D. A. (2007). Two faces of search: Alternative generation and alternative evaluation. \textit{Organization Science, 18}(1), 39--54. \doi{10.1287/orsc.1060.0216}

\bibitem[Kolko(2010)]{kolko2010} Kolko, J. (2010). Abductive thinking and sensemaking: The drivers of design synthesis. \textit{Design Issues, 26}(1), 15--28. \doi{10.1162/desi.2010.26.1.15}

\bibitem[Krippendorff(2018)]{krippendorff2018} Krippendorff, K. (2018). \textit{Content Analysis: An Introduction to Its Methodology} (4th ed.). SAGE.

\bibitem[Laumanns et~al.(2002)]{laumanns2002} Laumanns, M., Thiele, L., Deb, K., \& Zitzler, E. (2002). Combining convergence and diversity in evolutionary multiobjective optimization. \textit{Evolutionary Computation, 10}(3), 263--282. \doi{10.1162/106365602760234108}

\bibitem[Levinthal \& March(1993)]{levinthalmarch1993} Levinthal, D. A., \& March, J. G. (1993). The myopia of learning. \textit{Strategic Management Journal, 14}(S2), 95--112. \doi{10.1002/smj.4250141009}

\bibitem[Levinthal(1997)]{levinthal1997} Levinthal, D. A. (1997). Adaptation on rugged landscapes. \textit{Management Science, 43}(7), 934--950. \doi{10.1287/mnsc.43.7.934}

\bibitem[Liedtka(2015)]{liedtka2015} Liedtka, J. (2015). Perspective: Linking design thinking with innovation outcomes through cognitive bias reduction. \textit{Journal of Product Innovation Management, 32}(6), 925--938. \doi{10.1111/jpim.12163}

\bibitem[Loch et~al.(2001)]{loch2001} Loch, C. H., Terwiesch, C., \& Thomke, S. (2001). Parallel and sequential testing of design alternatives. \textit{Management Science, 47}(5), 663--678. \doi{10.1287/mnsc.47.5.663.10480}

\bibitem[Maher \& Poon(1996)]{maherpoon1996} Maher, M. L., \& Poon, J. (1996). Modeling design exploration as co-evolution. \textit{Microcomputers in Civil Engineering, 11}(3), 195--209. \doi{10.1111/j.1467-8667.1996.tb00323.x}

\bibitem[Mannor \& Tsitsiklis(2004)]{mannor2004} Mannor, S., \& Tsitsiklis, J. N. (2004). The sample complexity of exploration in the multi-armed bandit problem. \textit{Journal of Machine Learning Research, 5}, 623--648.

\bibitem[March(1991)]{march1991} March, J. G. (1991). Exploration and exploitation in organizational learning. \textit{Organization Science, 2}(1), 71--87. \doi{10.1287/orsc.2.1.71}

\bibitem[McComb et~al.(2015)]{mccomb2015} McComb, C., Cagan, J., \& Kotovsky, K. (2015). Lifting the veil: Drawing insights about design teams from a cognitively-inspired computational model. \textit{Design Studies, 40}, 119--142. \doi{10.1016/j.destud.2015.06.005}

\bibitem[Micheli et~al.(2019)]{micheli2019} Micheli, P., Wilner, S. J. S., Bhatti, S. H., Mura, M., \& Beverland, M. B. (2019). Doing design thinking: Conceptual review, synthesis, and research agenda. \textit{Journal of Product Innovation Management, 36}(2), 124--148. \doi{10.1111/jpim.12466}

\bibitem[Mintrom \& Luetjens(2016)]{mintrom2016} Mintrom, M., \& Luetjens, J. (2016). Design thinking in policymaking processes: Opportunities and challenges. \textit{Australian Journal of Public Administration, 75}(3), 391--402. \doi{10.1111/1467-8500.12211}

\bibitem[Nickerson(1998)]{nickerson1998} Nickerson, R. S. (1998). Confirmation bias: A ubiquitous phenomenon in many guises. \textit{Review of General Psychology, 2}(2), 175--220. \doi{10.1037/1089-2680.2.2.175}

\bibitem[Paton \& Dorst(2011)]{paton2011} Paton, B., \& Dorst, K. (2011). Briefing and reframing: A situated practice. \textit{Design Studies, 32}(6), 573--587. \doi{10.1016/j.destud.2011.07.002}

\bibitem[Rietzschel et~al.(2006)]{rietzschel2006} Rietzschel, E. F., Nijstad, B. A., \& Stroebe, W. (2006). Productivity is not enough: A comparison of interactive and nominal brainstorming groups on idea generation and selection. \textit{Journal of Experimental Social Psychology, 42}(2), 244--251. \doi{10.1016/j.jesp.2005.04.005}

\bibitem[Rittel \& Webber(1973)]{rittel1973} Rittel, H. W. J., \& Webber, M. M. (1973). Dilemmas in a general theory of planning. \textit{Policy Sciences, 4}(2), 155--169. \doi{10.1007/BF01405730}

\bibitem[Rudolph(1994)]{rudolph1994} Rudolph, G. (1994). Convergence analysis of canonical genetic algorithms. \textit{IEEE Transactions on Neural Networks, 5}(1), 96--101. \doi{10.1109/72.265964}

\bibitem[Russo et~al.(2018)]{russo2018} Russo, D. J., Van Roy, B., Kazerouni, A., Osband, I., \& Wen, Z. (2018). A tutorial on Thompson sampling. \textit{Foundations and Trends in Machine Learning, 11}(1), 1--96. \doi{10.1561/2200000070}

\bibitem[Sch{\"o}n(1983)]{schon1983} Sch{\"o}n, D. A. (1983). \textit{The Reflective Practitioner: How Professionals Think in Action}. Basic Books.

\bibitem[Seidel \& Fixson(2013)]{seidel2013} Seidel, V. P., \& Fixson, S. K. (2013). Adopting design thinking in novice multidisciplinary teams: The application and limits of design methods and reflexive practices. \textit{Journal of Product Innovation Management, 30}(S1), 19--33. \doi{10.1111/jpim.12061}

\bibitem[Shahriari et~al.(2016)]{shahriari2016} Shahriari, B., Swersky, K., Wang, Z., Adams, R. P., \& de Freitas, N. (2016). Taking the human out of the loop: A review of Bayesian optimization. \textit{Proceedings of the IEEE, 104}(1), 148--175. \doi{10.1109/JPROC.2015.2494218}

\bibitem[Simon(1956)]{simon1956} Simon, H. A. (1956). Rational choice and the structure of the environment. \textit{Psychological Review, 63}(2), 129--138. \doi{10.1037/h0042769}

\bibitem[Simon(1996)]{simon1996} Simon, H. A. (1996). \textit{The Sciences of the Artificial} (3rd ed.). MIT Press.

\bibitem[Smith \& Eppinger(1997)]{smitheppinger1997} Smith, R. P., \& Eppinger, S. D. (1997). Identifying controlling features of engineering design iteration. \textit{Management Science, 43}(3), 276--293. \doi{10.1287/mnsc.43.3.276}

\bibitem[Smith \& Winkler(2006)]{smithwinkler2006} Smith, J. E., \& Winkler, R. L. (2006). The optimizer's curse: Skepticism and postdecision surprise in decision analysis. \textit{Management Science, 52}(3), 311--322. \doi{10.1287/mnsc.1050.0451}

\bibitem[Sommer \& Loch(2004)]{sommerloch2004} Sommer, S. C., \& Loch, C. H. (2004). Selectionism and learning in projects with complexity and unforeseeable uncertainty. \textit{Management Science, 50}(10), 1334--1347. \doi{10.1287/mnsc.1040.0274}

\bibitem[Srinivas et~al.(2012)]{srinivas2012} Srinivas, N., Krause, A., Kakade, S. M., \& Seeger, M. W. (2012). Information-theoretic regret bounds for Gaussian process optimization in the bandit setting. \textit{IEEE Transactions on Information Theory, 58}(5), 3250--3265. \doi{10.1109/TIT.2011.2182033}

\bibitem[Suh(1990)]{suh1990} Suh, N. P. (1990). \textit{The Principles of Design}. Oxford University Press.

\bibitem[Sung \& Kelley(2019)]{sung2019} Sung, E., \& Kelley, T. R. (2019). Identifying design process patterns: A sequential analysis study of design thinking. \textit{International Journal of Technology and Design Education, 29}(2), 283--302. \doi{10.1007/s10798-018-9448-1}

\bibitem[Thomke(1998)]{thomke1998} Thomke, S. H. (1998). Managing experimentation in the design of new products. \textit{Management Science, 44}(6), 743--762. \doi{10.1287/mnsc.44.6.743}

\bibitem[Weitzman(1979)]{weitzman1979} Weitzman, M. L. (1979). Optimal search for the best alternative. \textit{Econometrica, 47}(3), 641--654. \doi{10.2307/1910412}

\bibitem[Woodbury \& Burrow(2006)]{woodbury2006} Woodbury, R. F., \& Burrow, A. L. (2006). Whither design space? \textit{Artificial Intelligence for Engineering Design, Analysis and Manufacturing, 20}(2), 63--82. \doi{10.1017/S0890060406060057}

\end{thebibliography}

\end{document}
